% GENERATED FILE. Edit canonical lessons, then run npm run generate:books.
% canonical-source-hash: fnv1a64:23a611c63f3c63e1
% canonical-lessons: JA-C498-kyuujitsu, JA-C498-fuyuyasumi, JA-C498-hirune, JA-C498-inemuri, JA-C498-yuukyuu

\chapter{A Day Off}
\label{ch:ja-a-day-off}
\hlchaptermodality{498}

\begin{chapteropening}
\textbf{By the end of this chapter:} \emph{I can talk about rest and holidays.}

Complete the last lesson of chapter 498: I can talk about rest and holidays.
\end{chapteropening}

\section[\texorpdfstring{\ja{きゅうじつ} (kyūjitsu)}{きゅうじつ (kyūjitsu)}]{\ja{きゅうじつ} (kyūjitsu) — a day off; a holiday}

\label{lesson:JA-C498-kyuujitsu}



\begin{quote}
Before the new one: say the Japanese for to roll into a ball, then the Japanese for a side dish (eaten with rice).
\end{quote}

\subsection*{You'll want to know: \ja{きゅうじつ}}
\textbf{\ja{きゅうじつ}} — \emph{kyūjitsu} — \textquotedblleft{}a day off; a holiday\textquotedblright{}.

\begin{quote}
\textbf{\ja{きゅうじつは} \ja{うちで} \ja{やすみます。}} — \emph{kyūjitsu wa uchi de yasumimasu} — \textquotedblleft{}On my days off I rest at home.\textquotedblright{}
\end{quote}

\begin{cousinweb}
\textbf{\ja{きゅうじつ}} (\emph{kyūjitsu}) is Sino-Japanese (kango): two Chinese characters read with their on'yomi, readings borrowed from Chinese. They are \emph{kyū}, \textquotedblleft{}rest\textquotedblright{} (Middle Chinese \emph{xjuw}) and \emph{jitsu}, \textquotedblleft{}day\textquotedblright{} (\emph{nyit}) (KANJIDIC2; Middle Chinese in Baxter's transcription).
\end{cousinweb}

\begin{grammarlens}[title={What you've built}]
One more word for describing what you lived and felt.
\end{grammarlens}

\subsection*{Your turn}
\begin{itemize}
\raggedright
  \item \emph{Say it:} \emph{kyūjitsu}
  \item \emph{Say it:} \emph{kyūjitsu}, once more
  \item \emph{Say it:} \emph{kyūjitsu wa uchi de yasumimasu} — On my days off I rest at home.
  \item \emph{Recall:} say the Japanese for to roll into a ball, then the Japanese for a side dish (eaten with rice), then say \emph{kyūjitsu} again
\end{itemize}

\begin{tcolorbox}[breakable,colback=teal!4,colframe=teal!35!black,title={Before you move on}]
What does \ja{きゅうじつ} mean? (\textquotedblleft{}A day off; a holiday\textquotedblright{}.) Say it once more.
\end{tcolorbox}
\section[\texorpdfstring{\ja{ふゆやすみ} (fuyuyasumi)}{ふゆやすみ (fuyuyasumi)}]{\ja{ふゆやすみ} (fuyuyasumi) — the winter holidays}

\label{lesson:JA-C498-fuyuyasumi}



\begin{quote}
Before the new one: say the Japanese for a side dish (eaten with rice), then the Japanese for a day off.
\end{quote}

\subsection*{You'll want to know: \ja{ふゆやすみ}}
\textbf{\ja{ふゆやすみ}} — \emph{fuyuyasumi} — \textquotedblleft{}the winter holidays\textquotedblright{}.

\begin{quote}
\textbf{\ja{ふゆやすみに} \ja{そふの} \ja{うちへ} \ja{いきます。}} — \emph{fuyuyasumi ni sofu no uchi e ikimasu} — \textquotedblleft{}In the winter holidays I go to my grandfather's house.\textquotedblright{}
\end{quote}

\begin{cousinweb}
\textbf{\ja{ふゆやすみ}} (\emph{fuyuyasumi}) is a native Japanese compound (wago), inherited rather than borrowed from Chinese: built from \emph{fuyu}, \textquotedblleft{}winter\textquotedblright{} and \emph{yasumi}, \textquotedblleft{}rest\textquotedblright{} (JMdict, KANJIDIC2).
\end{cousinweb}

\begin{grammarlens}[title={What you've built}]
One more word for describing what you lived and felt.
\end{grammarlens}

\subsection*{Your turn}
\begin{itemize}
\raggedright
  \item \emph{Say it:} \emph{fuyuyasumi}
  \item \emph{Say it:} \emph{fuyuyasumi}, once more
  \item \emph{Say it:} \emph{fuyuyasumi ni sofu no uchi e ikimasu} — In the winter holidays I go to my grandfather's house.
  \item \emph{Recall:} say the Japanese for a side dish (eaten with rice), then the Japanese for a day off, then say \emph{fuyuyasumi} again
\end{itemize}

\begin{tcolorbox}[breakable,colback=teal!4,colframe=teal!35!black,title={Before you move on}]
What does \ja{ふゆやすみ} mean? (\textquotedblleft{}The winter holidays\textquotedblright{}.) Say it once more.
\end{tcolorbox}
\section[\texorpdfstring{\ja{ひるね} (hirune)}{ひるね (hirune)}]{\ja{ひるね} (hirune) — a nap}

\label{lesson:JA-C498-hirune}



\begin{quote}
Before the new one: say the Japanese for a day off, then the Japanese for the winter holidays.
\end{quote}

\subsection*{You'll want to know: \ja{ひるね}}
\textbf{\ja{ひるね}} — \emph{hirune} — \textquotedblleft{}a nap\textquotedblright{}.

\begin{quote}
\textbf{\ja{ごはんの} \ja{あとで} \ja{ひるねを} \ja{します。}} — \emph{gohan no ato de hirune o shimasu} — \textquotedblleft{}I take a nap after eating.\textquotedblright{}
\end{quote}

\begin{cousinweb}
\textbf{\ja{ひるね}} (\emph{hirune}) is a native Japanese compound (wago), inherited rather than borrowed from Chinese: built from \emph{hiru}, \textquotedblleft{}noon\textquotedblright{} and \emph{ne}, \textquotedblleft{}sleep\textquotedblright{} (JMdict, KANJIDIC2).
\end{cousinweb}

\begin{grammarlens}[title={What you've built}]
One more word for describing what you lived and felt.
\end{grammarlens}

\subsection*{Your turn}
\begin{itemize}
\raggedright
  \item \emph{Say it:} \emph{hirune}
  \item \emph{Say it:} \emph{hirune}, once more
  \item \emph{Say it:} \emph{gohan no ato de hirune o shimasu} — I take a nap after eating.
  \item \emph{Recall:} say the Japanese for a day off, then the Japanese for the winter holidays, then say \emph{hirune} again
\end{itemize}

\begin{tcolorbox}[breakable,colback=teal!4,colframe=teal!35!black,title={Before you move on}]
What does \ja{ひるね} mean? (\textquotedblleft{}A nap\textquotedblright{}.) Say it once more.
\end{tcolorbox}
\section[\texorpdfstring{\ja{いねむり} (inemuri)}{いねむり (inemuri)}]{\ja{いねむり} (inemuri) — nodding off; dozing (while sitting)}

\label{lesson:JA-C498-inemuri}



\begin{quote}
Before the new one: say the Japanese for the winter holidays, then the Japanese for a nap.
\end{quote}

\subsection*{You'll want to know: \ja{いねむり}}
\textbf{\ja{いねむり}} — \emph{inemuri} — \textquotedblleft{}nodding off; dozing (while sitting)\textquotedblright{}.

\begin{quote}
\textbf{\ja{でんしゃで} \ja{いねむりを} \ja{しました。}} — \emph{densha de inemuri o shimashita} — \textquotedblleft{}I dozed off on the train.\textquotedblright{}
\end{quote}

\begin{cousinweb}
\textbf{\ja{いねむり}} (\emph{inemuri}) is a native Japanese compound (wago), inherited rather than borrowed from Chinese: built from \emph{i}, \textquotedblleft{}being, sitting\textquotedblright{} and \emph{nemuri}, \textquotedblleft{}sleep\textquotedblright{} (JMdict, KANJIDIC2).
\end{cousinweb}

\begin{grammarlens}[title={What you've built}]
One more word for describing what you lived and felt.
\end{grammarlens}

\subsection*{Your turn}
\begin{itemize}
\raggedright
  \item \emph{Say it:} \emph{inemuri}
  \item \emph{Say it:} \emph{inemuri}, once more
  \item \emph{Say it:} \emph{densha de inemuri o shimashita} — I dozed off on the train.
  \item \emph{Recall:} say the Japanese for the winter holidays, then the Japanese for a nap, then say \emph{inemuri} again
\end{itemize}

\begin{tcolorbox}[breakable,colback=teal!4,colframe=teal!35!black,title={Before you move on}]
What does \ja{いねむり} mean? (\textquotedblleft{}Nodding off; dozing (while sitting)\textquotedblright{}.) Say it once more.
\end{tcolorbox}
\section[\texorpdfstring{\ja{ゆうきゅう} (yūkyū)}{ゆうきゅう (yūkyū)}]{\ja{ゆうきゅう} (yūkyū) — paid leave; a paid holiday}

\label{lesson:JA-C498-yuukyuu}



\begin{quote}
Before the new one: say the Japanese for a nap, then the Japanese for nodding off.
\end{quote}

\subsection*{You'll want to know: \ja{ゆうきゅう}}
\textbf{\ja{ゆうきゅう}} — \emph{yūkyū} — \textquotedblleft{}paid leave; a paid holiday\textquotedblright{}.

\begin{quote}
\textbf{\ja{なつの} \ja{ゆうきゅうに} \ja{りょこうします。}} — \emph{natsu no yūkyū ni ryokō shimasu} — \textquotedblleft{}I go travelling on my summer paid leave.\textquotedblright{}
\end{quote}

\begin{cousinweb}
\textbf{\ja{ゆうきゅう}} (\emph{yūkyū}) is Sino-Japanese (kango): two Chinese characters read with their on'yomi, readings borrowed from Chinese. They are \emph{yū}, \textquotedblleft{}have\textquotedblright{} (Middle Chinese \emph{hjuwX}) and \emph{kyū}, \textquotedblleft{}salary\textquotedblright{} (\emph{kip}); JMdict gives \textquotedblleft{}paid\textquotedblright{} and \textquotedblleft{}paid leave\textquotedblright{} among its senses and marks it as a shortened word (JMdict, KANJIDIC2; Middle Chinese in Baxter's transcription).
\end{cousinweb}

\begin{grammarlens}[title={What you've built}]
One more word for describing what you lived and felt.
\end{grammarlens}

\subsection*{Your turn}
\begin{itemize}
\raggedright
  \item \emph{Say it:} \emph{yūkyū}
  \item \emph{Say it:} \emph{yūkyū}, once more
  \item \emph{Say it:} \emph{natsu no yūkyū ni ryokō shimasu} — I go travelling on my summer paid leave.
  \item \emph{Recall:} say the Japanese for a nap, then the Japanese for nodding off, then say \emph{yūkyū} again
\end{itemize}

\begin{tcolorbox}[breakable,colback=teal!4,colframe=teal!35!black,title={Before you move on}]
What does \ja{ゆうきゅう} mean? (\textquotedblleft{}Paid leave; a paid holiday\textquotedblright{}.) Say it once more.
\end{tcolorbox}
